\subsection{Backend mit FastAPI}
\label{subsec:fastapi-backend}

Das Backend bildet eine separate HTTP-Schicht für Web-cloud, Desktop-cloud und
Full-platform. FastAPI stellt unter \texttt{/api} derzeit nur
\texttt{/health} und \texttt{/ready} bereit. Produktspezifische Dienste,
Fachmodelle und Authentifizierung sind nicht vorgegeben
\parencite{fastapi2026features,kleiveist2026architecture}.

Liveness meldet den Prozesszustand. Bei aktiver Datenbankfähigkeit prüft
Readiness zusätzlich die Verbindung und antwortet bei Fehlern mit HTTP~503,
ohne interne Details offenzulegen. Router, Einstellungen und Abhängigkeit sind
getrennt testbar; Service- und Domänenmodule ergänzt erst das Produkt.

Im Deployment bleiben die Einheiten getrennt
(Abbildung~\ref{fig:deployment-architecture}): Der Vite-Build ist statisch
bereitstellbar, FastAPI läuft in einem eigenen nicht privilegierten Container.
Compose bildet diese providerneutrale Grenze lokal ab; PostgreSQL kommt nur
explizit hinzu
\parencite{kleiveist2026deployment}.

\begin{figure}[H]
  \centering
  \resizebox{0.96\textwidth}{!}{%
  \begin{tikzpicture}[node distance=8mm and 28mm]
    \node[casePrimary,minimum width=38mm] (vitesource) {Frontend-Quellen};
    \node[casePrimary,minimum width=38mm,right=of vitesource] (apisource) {FastAPI-Quellen};

    \node[caseBox,minimum width=46mm,below=of vitesource] (frontend)
      {\texttt{frontend/dist}\\Static Host oder Frontend-Image};
    \node[caseBox,minimum width=46mm,below=of apisource] (backend)
      {Backend-Image\\Uvicorn / FastAPI};

    \node[caseBox,minimum width=34mm,below left=of backend] (health)
      {\texttt{/api/health}\\prozessbezogen};
    \node[caseBox,minimum width=34mm,below=of backend] (readiness)
      {\texttt{/api/ready}\\abhängigkeitsbezogen};
    \node[caseOptional,minimum width=38mm,below=of readiness] (postgres)
      {PostgreSQL\\explizites Compose-Profil};

    \draw[caseFlow] (vitesource) -- node[left,font=\sffamily\tiny] {Vite-Build} (frontend);
    \draw[caseFlow] (apisource) -- node[right,font=\sffamily\tiny] {Container-Build} (backend);
    \draw[caseFlow] (frontend) -- node[above,font=\sffamily\tiny] {HTTP /api} (backend);
    \draw[caseSecondaryFlow] (backend) -- (health);
    \draw[caseSecondaryFlow] (backend) -- (readiness);
    \draw[caseSecondaryFlow] (readiness) -- node[right,font=\sffamily\tiny] {SELECT 1, falls aktiv} (postgres);

    \begin{scope}[on background layer]
      \node[caseGroup,fit=(frontend)(backend)(health)(readiness)(postgres),
        label={[font=\sffamily\bfseries\scriptsize,text=caseBlue]below:providerneutrale Einheiten; Docker/Compose als lokale Betriebsgrenze}] {};
    \end{scope}
  \end{tikzpicture}}
  \caption{Providerneutrale Deployment-Architektur}
  \label{fig:deployment-architecture}
  \smallskip
  {\footnotesize Quelle: Eigene Darstellung auf Grundlage von
  \textcite{kleiveist2026deployment}.\par}
\end{figure}

\beginnerinput{workspace/statement/300_architektur/330}
